\IfFileExists{papers/formalized-vs-literature-preamble.tex}{\documentclass[11pt]{article}
\usepackage[utf8]{inputenc}
\usepackage[T1]{fontenc}
\usepackage{amsmath,amssymb,mathtools}
\usepackage{booktabs,tabularx,array}
\usepackage{enumitem}
\usepackage[margin=1in]{geometry}
\usepackage{xcolor}
\usepackage[hidelinks]{hyperref}
\usepackage{microtype}
\setlength{\parindent}{0pt}
\setlength{\parskip}{0.55em}
\newcommand{\Lean}[1]{\texttt{\nolinkurl{#1}}}
\newcommand{\RepoFile}[1]{\texttt{\nolinkurl{#1}}}
\newcommand{\Class}[1]{\textbf{#1}}
\newcommand{\Top}{\mathord{\top}}
\newcommand{\R}{\mathbb{R}}
\newcommand{\ENN}{\mathbb{R}_{\ge 0}^{\infty}}
\newcommand{\KL}{\operatorname{KL}}
\newcommand{\W}{\mathsf{W}}
\newcommand{\statusline}[2]{%
  \begin{center}\fbox{\begin{minipage}{0.92\linewidth}\textbf{Completion class:} #1\\[2pt]
  \textbf{Strongest caveat:} #2\end{minipage}}\end{center}}
\newenvironment{comparison}{\tabularx{\linewidth}{>{\bfseries}p{0.25\linewidth}X}}{\endtabularx}
\newcommand{\snapshot}{\texttt{902f51aa01a737af68f6feb71c83263cc9b1f4d9}}
}{\documentclass[11pt]{article}
\usepackage[utf8]{inputenc}
\usepackage[T1]{fontenc}
\usepackage{amsmath,amssymb,mathtools}
\usepackage{booktabs,tabularx,array}
\usepackage{enumitem}
\usepackage[margin=1in]{geometry}
\usepackage{xcolor}
\usepackage[hidelinks]{hyperref}
\usepackage{microtype}
\setlength{\parindent}{0pt}
\setlength{\parskip}{0.55em}
\newcommand{\Lean}[1]{\texttt{\nolinkurl{#1}}}
\newcommand{\RepoFile}[1]{\texttt{\nolinkurl{#1}}}
\newcommand{\Class}[1]{\textbf{#1}}
\newcommand{\Top}{\mathord{\top}}
\newcommand{\R}{\mathbb{R}}
\newcommand{\ENN}{\mathbb{R}_{\ge 0}^{\infty}}
\newcommand{\KL}{\operatorname{KL}}
\newcommand{\W}{\mathsf{W}}
\newcommand{\statusline}[2]{%
  \begin{center}\fbox{\begin{minipage}{0.92\linewidth}\textbf{Completion class:} #1\\[2pt]
  \textbf{Strongest caveat:} #2\end{minipage}}\end{center}}
\newenvironment{comparison}{\tabularx{\linewidth}{>{\bfseries}p{0.25\linewidth}X}}{\endtabularx}
\newcommand{\snapshot}{\texttt{902f51aa01a737af68f6feb71c83263cc9b1f4d9}}
}
\title{Positive matrix scaling and discrete Schrödinger potentials: formalized result versus the literature}
\author{}
\date{Formalization snapshot \snapshot}
\begin{document}
\maketitle
\statusline{\Class{G/S}: complete finite strictly-positive rectangular existence theorem and square Schrödinger-potential specialization}{The theorem does not cover nonnegative matrices with support conditions, and its proof is an independent log-domain convex-minimization route rather than Sinkhorn--Knopp convergence.}
\section{Literature targets}
The relevant literature includes Sinkhorn's scaling theorem, Sinkhorn--Knopp (1967), Franklin--Lorenz (1989), and the finite Schrödinger system used by Chen--Georgiou--Pavon.  The local source reconstructions are
\RepoFile{prose/distilled_literature/SinkhornKnopp1967_matrix_scaling.tex},
\RepoFile{prose/distilled_literature/FranklinLorenz1989_matrix_scaling.tex}, and
\RepoFile{prose/distilled_literature/CGP2021_sinkhorn_schrodinger_bridge.tex}.

\section{Lean declarations}
\begin{itemize}[leftmargin=2em]
\item \Lean{ForMathlib.matrix_scaling_exists_rectangular};
\item \Lean{ForMathlib.matrix_scaling_exists};
\item \Lean{ForMathlib.sinkhorn_potentials_exist};
\item \Lean{ChenGeorgiouPavon2021.sinkhorn_potentials_exist}.
\end{itemize}
The core is in \RepoFile{ForMathlib/LinearAlgebra/Matrix/SinkhornScaling.lean}.

\section{What Lean proves}
For finite row and column index types, strictly positive prescribed marginals with equal total mass, and a strictly positive rectangular kernel $G$, Lean constructs positive scalings $a,b$ such that
\[
 a_i\sum_jG_{ij}b_j=p_i,
 \qquad
 b_j\sum_iG_{ij}a_i=q_j.
\]
The square theorem and four-potential Schrödinger-system theorem are thin specializations.

\section{Proof architecture}
The proof minimizes a log-domain convex potential on an open simplex.  Boundary blow-up confines a sublevel set to a compact region; the extreme-value theorem supplies a minimizer; one-dimensional directional derivatives give the marginal equations; equal total mass removes the Lagrange multiplier.  This is not a proof by convergence of alternating row/column normalization.

\section{Comparison}
\begin{comparison}
Matrix entries & Strictly positive only.\\
Shape & Rectangular core and square wrapper.\\
Marginals & Strictly positive with equal mass.\\
Existence & Complete.\\
Uniqueness up to gauge & Proved elsewhere for finite potential systems, not bundled into this existence theorem.\\
Algorithmic convergence & Separate theorem family.\\
Support-sensitive case & Not formalized.\\
\end{comparison}

\section{Nonclaims}
No claim is made for matrices with zeros, total support, full indecomposability, or the exact support criteria of Sinkhorn--Knopp.  The proof also does not show that raw Sinkhorn iterates converge to the constructed scaling.
\end{document}
